\newcommand{\documento}{\LdP}

%Generali
\newcommand{\progetto}{QuizziPedia}
\newcommand{\gruppo}{TheFellowshipOfTheCode}
\newcommand{\gruppoLink}{\href{http://thefellowshipofthecode.github.io/}{TheFellowshipOfTheCode}}
\newcommand{\email}{\href{mailto:thefellowshipofthecode@gmail.com}{thefellowshipofthecode@gmail.com}}

%Documenti
\newcommand{\AdR}{Analisi dei Requisiti}
\newcommand{\NdP}{Norme di Progetto}
\newcommand{\PdP}{Piano di Progetto}
\newcommand{\SdF}{Studio di Fattibilità}
\newcommand{\PdQ}{Piano di Qualifica}
\newcommand{\VI}{Verbale Interno}
\newcommand{\VE}{Verbale Esterno}
\newcommand{\ST}{Specifica Tecnica}
\newcommand{\DDP}{Definizione di Prodotto}
\newcommand{\MU}{Manuale Utente}
\newcommand{\G}{Glossario}
\newcommand{\LdP}{Lettera di Presentazione}
\newcommand{\NdPv}{NormeDiProgetto\_v\_1\_0\_0}
\newcommand{\PdPv}{PianoDiProgetto\_v\_1\_0\_0}
\newcommand{\PdQv}{PianoDiQualifica\_v\_1\_0\_0}
\newcommand{\SdFv}{StudioDiFattibilità\_v\_1\_0\_0}

%Componenti del gruppo
\newcommand{\AF}{Alberto Ferrara}
\newcommand{\SM}{Simone Magagna}
\newcommand{\FB}{Franco Berton}
\newcommand{\MP}{Marco Prelaz}
\newcommand{\MV}{Mattia Varotto}
\newcommand{\GN}{Matteo Gnoato}
\newcommand{\GR}{Matteo Granzotto}

%Ruoli
\newcommand{\RdP}{Responsabile di Progetto}
\newcommand{\Res}{Responsabile}
\newcommand{\Amm}{Amministratore}
\newcommand{\Ver}{Verificatore}
\newcommand{\Prog}{Progettista}
\newcommand{\Progr}{Programmatore}
\newcommand{\Ana}{Analista}
\newcommand{\RdPs}{Responsabili di Progetto}
\newcommand{\Ress}{Responsabile}
\newcommand{\Amms}{Amministratori}
\newcommand{\Vers}{Verificatori}
\newcommand{\Progs}{Progettisti}
\newcommand{\Progrs}{Programmatori}
\newcommand{\Anas}{Analisti}

%Professori e proponente
\newcommand{\TV}{Prof. Tullio Vardanega}
\newcommand{\RC}{Prof. Riccardo Cardin}
\newcommand{\ZU}{Zucchetti S.P.A.}
\newcommand{\proponente}{Zucchetti S.P.A.}

\newcommand{\diaryEntry}[5]{#2 & \emph{#4} & #3 & #5 & #1\\ \hline}

%comando per una nuova riga nella tabella del diario delle modifiche
\newcommand{\specialcell}[2][c]{%
	\begin{tabular}[#1]{@{}c@{}}#2\end{tabular}}

\renewcommand*\sectionmark[1]{\markboth{#1}{}}
\renewcommand*\subsectionmark[1]{\markright{#1}}

%Pediodi di lavoro 
\newcommand{\AR}{Analisi dei Requisiti}
\newcommand{\AD}{Analisi dei Requisiti in Dettaglio}
\newcommand{\PA}{Progettazione Architetturale}
\newcommand{\PD}{Progettazione di Dettaglio}
\newcommand{\CO}{Codifica}
\newcommand{\VV}{Validazione}

% Revisioni
\newcommand{\RR}{Revisione dei Requisiti}
\newcommand{\RP}{Revisione di Progettazione}
\newcommand{\RQ}{Revisione di Qualifica}
\newcommand{\RA}{Revisione di Accettazione}

% Comandi analisi dei requisiti
\newcommand{\uau}{utente autenticato}
\newcommand{\uaus}{utenti autenticati}
\newcommand{\uaupro}{utente autenticato pro}
\newcommand{\uauspro}{utenti autenticati pro}

\newcommand{\myincludegraphics}[2][]{%
	\setbox0=\hbox{\phantom{X}}%
	\vtop{
		\hbox{\phantom{X}}
		\vskip-\ht0
		\hbox{\includegraphics[#1]{#2}}}}

\renewcommand\footnoterule{\rule{\linewidth}{1pt}}

%----------------------------------------------------------------------------------------
%	DOCUMENT CONFIGURATIONS
%----------------------------------------------------------------------------------------

\documentclass{letter}
\usepackage{eurosym}
\usepackage{fancyhdr}
\usepackage{graphicx}
\usepackage[utf8]{inputenc}
% Adjust margins for aesthetics
\addtolength{\voffset}{-0.5in}
\addtolength{\hoffset}{-0.3in}
\addtolength{\textheight}{2cm}

%\longindentation=0pt % Un-commenting this line will push the closing "Sincerely," to the left of the page
\fancypagestyle{firstpage}{
	% cancella tutti i campi di intestazione e piè di pagina
	\fancyhf{}
	\renewcommand{\headrulewidth}{0.0pt}
	\renewcommand{\footrulewidth}{0.3pt}

	\fancyhead[L]{
		\includegraphics[height=3cm, keepaspectratio=true]{../../Template/icone/LogoGruppo.png}
	}
	\fancyhead[R]{
		\Large{\textbf{Gruppo \\ \gruppo{}}} \\
		\vspace{55pt}
	}

	\cfoot{ \small{
			Gruppo: \gruppo{} \\ E-mail: TheFellowshipOfTheCode@gmail.com
		}}
	}

%----------------------------------------------------------------------------------------
%	YOUR NAME & ADDRESS SECTION
%----------------------------------------------------------------------------------------

	\signature{
		\begin{center}
			\textit{\SM{} \\ \Res \gruppo{}}
		\end{center}
	}


%----------------------------------------------------------------------------------------

\begin{document}

%----------------------------------------------------------------------------------------
%	ADDRESSEE SECTION
%----------------------------------------------------------------------------------------

\begin{letter}{\hfill \TV \\
	\hfill \RC \\
	\hfill Università degli Studi di Padova \\
	\hfill Dipartimento di Matematica Pura ed Applicata \\
	\hfill Via Trieste, 63 \\
	\hfill 35121 Padova \\
	\hfill 14 Giugno 2016}

%----------------------------------------------------------------------------------------
%	LETTER CONTENT SECTION
%----------------------------------------------------------------------------------------
\date{}
\opening{Egregio \TV,}

	\begin{quotation}
		Con la presente il gruppo \textit{\gruppo} intende comunicarLe ufficialmente la partecipazione alla \textbf{Revisione di Accettazione} per la realizzazione del prodotto da Lei commissionato, denominato \textbf{\progetto}.
		\\
		Si allegano alla presente lettera i seguenti documenti:
		\begin{itemize}
			\item \textit{StudioDiFattibilità\_v1\_0\_0};
			\item \textit{NormeDiProgetto\_v5\_0\_0};
			\item \textit{PianoDiProgetto\_v6\_0\_0.pdf};
			\item \textit{PianoDiQualifica\_v6\_0\_0.pdf};
			\item \textit{AnalisiDeiRequisiti\_v3\_0\_0};
			\item \textit{DefinizioneDiProdotto\_v4\_0\_0};
			\item \textit{Glossario\_v4\_0\_0.pdf};
			\item \textit{ManualeUtente\_v2\_0\_0.pdf};
			\item \textit{ManualeManutentore\_v2\_0\_0.pdf};
			\item \textit{Verbale\_1\_Interno\_2015-12-03.pdf};
			\item \textit{Verbale\_2\_Interno\_2015-12-11.pdf};
			\item \textit{Verbale\_3\_Interno\_2015-12-29.pdf};
			\item \textit{Verbale\_1\_Esterno\_2016-01-11.pdf};
			\item \textit{Verbale\_4\_Interno\_2016-01-12.pdf};
			\item \textit{Verbale\_5\_Interno\_2016-02-25.pdf};
			\item \textit{Verbale\_6\_Interno\_2016-03-01.pdf};
			\item \textit{Verbale\_7\_Interno\_2016-03-07.pdf};
			\item \textit{Verbale\_2\_Esterno\_2016-03-11.pdf};
			\item \textit{Verbale\_8\_Interno\_2016-03-21.pdf};
			\item \textit{Verbale\_9\_Interno\_2016-03-25.pdf};
			\item \textit{Verbale\_10\_Interno\_2016-04-26.pdf};
			\item \textit{Verbale\_11\_Interno\_2016-05-02.pdf};
			\item \textit{Verbale\_12\_Interno\_2016-05-09.pdf};	
			\item \textit{Verbale\_13\_Interno\_2016-05-27.pdf};
			\item \textit{Verbale\_14\_Interno\_2016-06-06.pdf}.
		\end{itemize}
		e il codice prodotto presente nella cartella \textit{Codice}.
		\newpage
		I documenti consegnati trattano in maniera approfondita delle norme di progetto, dei dettagli di qualifica, di requisiti e di prodotto.
		Il codice prodotto rappresenta l'implementazione dell'applicazione \progetto. Il costo finale del lavoro effettuato dal gruppo è di \euro13.468, come analizzato nel \textit{\PdP}.	
	\end{quotation}
	\closing{Cordiali saluti,}
	\includegraphics[keepaspectratio = true, width=6cm]{Firma/SM.png}

%----------------------------------------------------------------------------------------

\end{letter}

\end{document}
